\section{Machine-checked: always-embeddability}
\label{app:lean}

The declaration heads below summarize the Lean~4 development in
\nolinkurl{lean/always_embeddable.lean} that establishes \cref{thm:design:embed}
(there named \texttt{embedding\_closed\_under\_context}),
faithful to the abstract definitions of \cref{sec:design:abstract}.
The disjointness of openers and closers assumed there ($O \cap C = \emptyset$)
appears as the premise \texttt{hdisj};
tactic proof bodies are elided and given in full in the source file.

\begin{lstlisting}
import Mathlib

variable {α : Type*}         -- the alphabet
variable (Pi : Set (α × α))  -- the set of pairs

def Opener     (x : α) : Prop := ∃ y, (x, y) ∈ Pi
def Closer     (x : α) : Prop := ∃ y, (y, x) ∈ Pi
def Matcher    (x : α) : Prop := Opener Pi x ∨ Closer Pi x
def Nonmatcher (x : α) : Prop := ¬ Matcher Pi x

-- Lemma: no character is both an opener and a closer.
theorem disjoint (h : ∀ x y z, (x, y) ∈ Pi → (z, x) ∈ Pi → False)
    (x : α) : ¬ (Opener Pi x ∧ Closer Pi x) := by ...

-- The language of matchertext strings,
-- one constructor per clause of the inductive definition.
inductive MT (Pi : Set (α × α)) : List α → Prop where
  | flat (n : List α) (h : ∀ x ∈ n, Nonmatcher Pi x) : MT Pi n
  | nest (m₁ m₂ m₃ : List α) (o c : α) (hp : (o, c) ∈ Pi)
         (h₁ : MT Pi m₁) (h₂ : MT Pi m₂) (h₃ : MT Pi m₃) :
      MT Pi (m₁ ++ [o] ++ m₂ ++ [c] ++ m₃)

-- Lemma: the empty string is matchertext.
theorem mt_nil : MT Pi ([] : List α) := by ...

-- Lemma: the language is closed under concatenation.
theorem mt_append {a b : List α} (ha : MT Pi a) (hb : MT Pi b) :
    MT Pi (a ++ b) := by ...

-- Lemma: wrapping a matchertext string in a single matched pair.
theorem mt_wrap {m : List α} (hm : MT Pi m) {o c : α} (hp : (o, c) ∈ Pi) :
    MT Pi ([o] ++ m ++ [c]) := by ...

-- Matched delimiter contexts: nested matcher pairs
-- with nonmatchers interspersed freely on either side.
inductive AllMatch (Pi : Set (α × α)) :
    List α → List α → Prop where
  | nil : AllMatch Pi [] []
  | cons (o c : α) (s_o' s_c' : List α) (hp : (o, c) ∈ Pi)
         (rest : AllMatch Pi s_o' s_c') :
      AllMatch Pi (o :: s_o') (s_c' ++ [c])
  | nonmatch_o (x : α) (s_o' s_c' : List α) (hx : Nonmatcher Pi x)
         (rest : AllMatch Pi s_o' s_c') :
      AllMatch Pi (x :: s_o') s_c'
  | nonmatch_c (y : α) (s_o' s_c' : List α) (hy : Nonmatcher Pi y)
         (rest : AllMatch Pi s_o' s_c') :
      AllMatch Pi s_o' (s_c' ++ [y])

-- Theorem (always-embeddability): embedding a matchertext
-- string in a matched context always yields matchertext.
theorem embedding_closed_under_context
    (hdisj : ∀ x y z, (x, y) ∈ Pi → (z, x) ∈ Pi → False)
    {m : List α} (hm : MT Pi m)
    {s_o s_c : List α} (hs : AllMatch Pi s_o s_c) :
    MT Pi (s_o ++ m ++ s_c) := by ...
\end{lstlisting}

\paragraph{Proof sketches of the lemmas}

\begin{proof}[Proof sketch of \cref{lem:design:concat}]
By induction on the derivation of $a \in L$.
If $a$ is a string of nonmatchers,
a second induction on the derivation of $b \in L$
either extends that nonmatcher string, when $b$ is also flat,
or splits off the outermost matched pair of $b$
and reattaches $a$ on its left via the nesting rule.
If $a$ itself has the form $a_1||o||a_2||c||a_3$,
then $a||b = a_1||o||a_2||c||(a_3||b)$,
which lies in $L$ by the induction hypothesis on $a_3$
and one application of the nesting rule.
\end{proof}

\begin{proof}[Proof sketch of \cref{lem:design:wrap}]
Immediate from the nesting rule of \cref{sec:design:abstract}
with $m_1$ and $m_3$ both empty and $m_2 = m$.
\end{proof}
